\documentclass[12pt]{article}
\usepackage[utf8]{inputenc}
\usepackage[french]{babel}
\usepackage{amsmath}
\usepackage{pgfplots}
\usepackage{circuitikz}

\begin{document}

\section*{Circuit RC et Charge du Condensateur}

\subsection*{Schéma du Circuit RC}
\begin{center}
\begin{tikzpicture}
    % Schéma simplifié du circuit RC
    \draw (0,0) node[ground] {} -- (0,2) to[V, l=$\mathcal{V}_0$] (0,2) -- (3,2) to[R, l=$R$] (3,0) -- (0,0);
    \draw (3,2) to[C, l=$C$] (3,0);
\end{tikzpicture}
\end{center}

\subsection*{Graphique de la Charge}
\begin{center}
\begin{tikzpicture}
\begin{axis}[
    xlabel={Temps $t$ (s)},
    ylabel={Tension $V_C(t)$ (V)},
    xmin=0, xmax=0.01,
    ymin=0, ymax=10,
    samples=100,
    grid=both,
    width=10cm,
    height=6cm
]
% Courbe de charge : V(t) = V0 * (1 - exp(-t/RC))
\addplot[blue, thick] {10*(1 - exp(-x/0.005))};
\addlegendentry{$V_C(t) = V_0(1 - e^{-t/RC})$};
\end{axis}
\end{tikzpicture}
\end{center}

\subsection*{Équation Théorique}
La tension aux bornes du condensateur suit la loi suivante :
$$
V_C(t) = V_0 \left(1 - e^{-\frac{t}{RC}}\right)
$$

\end{document}